\section{Results}
\label{sec:results}

\newcommand{\maybetable}[1]{\IfFileExists{tables/#1.tex}{\input{tables/#1}}{\TODO{table #1}}}

\begin{table}[t]
\centering\small
\caption{Degradation families and their parameters at levels 1--3 (level 0 is the unmodified policy).
Latency is in replanning steps of $0.5$\,s; single-shot evaluators cap it at three steps.}
\label{tab:degradations}
\begin{tabular}{@{}llccc@{}}
\toprule
Family & Parameter & L1 & L2 & L3 \\
\midrule
Camera dropout & cameras blanked & 1 & 2 & 4 \\
Image noise & $\sigma$ (pixel value) & 10 & 25 & 50 \\
Stale observation & age (s) & 0.5 & 1.0 & 1.5 \\
Lateral noise & peak drift (m) & 0.3 & 0.8 & 1.5 \\
Late braking & path stretch & 0.15 & 0.40 & 0.70 \\
Latency & steps & 1 & 2 & 4 (NAVSIM: 3) \\
\midrule
Tail lateral (benign) & peak drift (m) & 0.3 & 0.8 & 1.5 \\
Tail stretch (benign) & path stretch & 0.15 & 0.40 & 0.70 \\
\bottomrule
\end{tabular}
\end{table}

\subsection{Which degradations give a ground truth?}
\label{sec:results:mono}
\maybetable{monotonicity}
\TODO{Text from the split analysis: retained families on the selection half, and the all-scenario
analysis as a sensitivity check.}

\subsection{Do evaluators recover the known order?}
\label{sec:results:order}
\maybetable{ordering}
\TODO{Ordering recovery and sample efficiency per evaluator on the evaluation half; H1.}

\subsection{Specificity}
\label{sec:results:spec}
\maybetable{specificity}
\TODO{Benign validation on the reference; verdicts per evaluator.}

\subsection{Ranking real policies}
\label{sec:results:popB}
\maybetable{popB}
\TODO{Family contrasts on the shared scenes; agreement; per-scenario correlation against the
test-retest ceiling.}
